% ------------------------------------------------------------
% Page 113
% ------------------------------------------------------------

« Pour ses goûts d’agriculteur, Ferrussac était un nouveau théâtre où s’exercer : la nature
schisteuse du sol appelait des amendements calcaires (1), des prairies naturellement irriguées
demandaient des soins bien autres que les pacages secs ou les champs à céréales du Causse,
les flancs des montagnes se prêtaient au reboisement par des arbres verts.
L’habitation, elle-même, pour être adaptée à la vie simple mais confortable des hôtes appelait
des améliorations fondamentales et des accessoires d’agrément. Il semblait devoir vivre
longtemps tant il paraissait solide… »

Mme. Richard-Valdeyron et M. Jacky Martin dont les arrière-grands-parents habitaient Rousses puis Ferrussac
m’ont raconté quelques souvenirs transmis par leurs grands-parents

\begin{tcolorbox}
Mr. Cambessédès disposait de domestiques dont un noir. Toutes les terres de Ferrussac lui appartenaient et
plusieurs fermiers les mettaient en valeur lui assurant des rentes suffisantes pour son train de vie qu’il avait
modeste. Il avait installé une balançoire dans un magnifique sapin. En fait il s’agissait d’un hamac placé entre
deux grands arbres qui agrémentaient le chemin d’accès au château. Non loin une grande volière abritait
toutes sortes d’oiseaux, exotiques pour certains. Un jour, en Août 1862 il dégringola de son hamac et sa santé
déclina.

Ses héritiers continuèrent pendant plus d’un demi-siècle à percevoir les fermages des différents exploitants. Mais
leur fortune périclita peut-être à la suite des guerres de 1870 et 1914, mais également, dit-on, parce qu’ils
s’adonnaient au jeu et, vers 1920 ils durent vendre.
Les fermiers du château purent racheter leurs fermes avec les terres.
« Mes arrière-grands-parents, les Valdeyron, fermiers à Aures sur le Causse Méjean acquirent ainsi le château. »
Les grands parents de Jacky, Numa Martin frère d’Elie de Raffègues, ceux d’Odette Brun-Bertrand et de
Jeannine Martin-Bertrand de Carnac (Rousses) Jules et Délie Bertrand, la famille Gache... Furent dès lors
propriétaires exploitants.
\end{tcolorbox}

A la suite de cette chute de hamac « sa santé fut atteinte à la source. Après un traitement
local, se succédèrent, érysipèle, asthme, et œdème qui l’emporta le 20 octobre 1863, précise
E. Planchon, entouré des tendres soins d’une femme et d’une sœur aimées, d’un neveu que
son affection avait fait, d’avance, son fils et son héritier. Il s’éteignit sans lutte et comme sans
souffrance dans la nuit du 20 Octobre 1863.

Après sa mort Mme. Cambessédès et Paul de Froment, son légataire universel, remirent à la
Faculté des Sciences de Montpellier, tous ses herbiers, une collection considérable de
nombreux pays, du Brésil aux Indes en passant par les Antilles ou Madagascar, que lui
avaient confiés pour étude de nombreux confrères botanistes.

\footnote{« Jacques Cambessédès a pu voir tout juste les premiers résultats de cette transformation du sol. Nous lui
avons apporté, écrit sa digne compagne, dans sa chambre de malade, le froment qu’il avait fait semer sous ses
yeux et qui venait d’être dépique. Amélioration immense pour un pays qui n’avait jamais produit que du seigle ».}
